\chapter*{Lời Tựa}
\addcontentsline{toc}{chapter}{Lời Tựa}
\markboth{Lời Tựa}{Lời Tựa}

\begin{truyen}{Lời Tựa}{One Hundred Human Paradoxes}
\chuhoa{C}{àng sống lâu, người ta càng hiểu đời không đi theo những câu trả lời một màu. Có cái được mà cũng là mất. Có cái thua mà lại là lối lớn lên. Có cái đúng về lý nhưng sai vì thiếu lòng người.}
\textit{(The longer we live, the more we see that life does not move in one-color answers. Some gains are also losses. Some defeats become ways of growing. Some things are right in logic but wrong because they lack human warmth.)}

Cuốn sách này mở thành ba phần lớn: \textbf{Được và Mất}, \textbf{Thắng và Thua}, \textbf{Đúng và Sai}. Từ đó, nó đi tiếp sang những nghịch lý khác của đời sống như \textbf{Có và Không}, \textbf{May và Rủi}, \textbf{Nhanh và Chậm}, \textbf{Mạnh và Yếu}, \textbf{Thực và Ảo}. Mỗi phần đều đi từ những lát cắt rất gần gũi của gia đình, trường học, tình cảm, công việc và lòng tự trọng.
\textit{(This book unfolds into three major parts: \textbf{Gain and Loss}, \textbf{Winning and Losing}, and \textbf{Right and Wrong}. From there it extends into other paradoxes of life such as \textbf{Having and Not Having}, \textbf{Luck and Misfortune}, \textbf{Fast and Slow}, \textbf{Strong and Weak}, and \textbf{Real and Illusion}. Each part grows out of familiar slices of family, school, love, work, and dignity.)}

Điều cuốn sách muốn giữ lại không phải là những câu châm ngôn cho đẹp, mà là một cách nhìn tỉnh hơn: cái được thường kéo theo cái mất, cái thắng có thể là cái thua, và cái đúng nhiều khi đòi hỏi phải trả giá.
\textit{(What this book wants to preserve is not decorative maxims, but a clearer way of seeing: gain often carries loss, victory can hide defeat, and what is right often demands a cost.)}
\end{truyen}

\begin{baihoc}
Sâu sắc không nằm ở chỗ đời có ít nghịch lý, mà ở chỗ con người chịu nhìn thẳng vào chúng để sống khôn hơn.
\textit{(Depth does not lie in life having few paradoxes, but in our willingness to face them so we can live more wisely.)}
\end{baihoc}

\begin{ghinhoanh}
\textbf{Short English to remember:}

Life is deeper than one-sided answers.

Paradoxes are where wisdom begins.
\end{ghinhoanh}

\begin{tuhocnhanh}
1. Trong ba cặp chủ đề này, em thấy mình đang mắc kẹt ở cặp nào nhất? \textit{(Among these three paradox pairs, where do you feel most stuck right now?)}

2. Sau khi đọc cuốn này, em muốn sống khác đi ở điểm nào rõ nhất? \textit{(After reading this book, what is the clearest way you want to live differently?)}
\end{tuhocnhanh}

\begin{flushright}
\textbf{Nguyễn Văn Sang}\\
\textit{GV THPT Nguyễn Hữu Cảnh - TP HCM}
\end{flushright}

\ngancach